\section{Related Work}

Prior network intrusion studies span classical one-class detectors, deep sequence models, and adversarial anomaly detectors. IsolationForest \cite{liu2008iforest} and OneClassSVM \cite{scholkopf2001ocsvm} remain common low-label baselines, while online anomaly detectors such as Kitsune emphasize deployment practicality under streaming constraints \cite{mirsky2018kitsune}. Deep baselines include GANomaly \cite{akcay2018ganomaly}, TranAD \cite{tuli2022tranad}, and related Transformer-style anomaly models such as Anomaly Transformer and TimesNet \cite{xu2022anomalytransformer,wu2023timesnet}.

Many of these studies are reported at flow-level or ranking-focused settings. This is informative for representation quality, but it does not automatically establish stable alert behavior under a fixed operating threshold. Our scope differs by prioritizing window-level alarm behavior and explicit benign-noise reporting under fixed-threshold deployment assumptions.

Deep anomaly-detection literature also includes one-class and reconstruction-style methods, such as DeepSVDD \cite{ruff2018deep}, AnoGAN \cite{schlegl2017anogan}, and f-AnoGAN \cite{schlegl2019fanogan}. These approaches improve representation learning under limited labels, but deployment quality still depends on threshold policy and benign alarm control, not ranking metrics alone.

For temporal modeling, TCNs provide efficient sequence processing \cite{bai2018tcn}. In adversarial settings, WGAN and WGAN-GP improve critic training stability compared with vanilla GAN objectives \cite{goodfellow2014gan,arjovsky2017wgan,gulrajani2017wgangp}. Our method follows this line but emphasizes calibrated alarm operation through benign-only independent threshold calibration, rather than architecture-only gains.

Flow-based and Transformer-based NIDS studies further show that model architecture choices can improve representation quality while leaving deployment threshold transfer as an open issue \cite{manocchio2024flowtransformer,long2024transformercloud,xi2024multiscale}. This motivates our scope: representative strict-protocol baselines with explicit calibration and observed benign FPR reporting.

Dataset-wise, CICIDS2017 is widely used for contemporary flow-based NIDS evaluation \cite{sharafaldin2018cicids}. SWaT and UNSW-NB15 provide additional but distinct validation contexts, including industrial-control telemetry and distribution-shift stress conditions \cite{goh2016swat,moustafa2015unsw}. In this paper, these supplementary datasets are used to discuss bounded applicability, not comprehensive cross-dataset generalization.

A related practical thread is threshold calibration and false-positive control in operational anomaly detection. One-class and online IDS methods typically require explicit threshold selection and post-deployment monitoring of benign alert rates \cite{scholkopf2001ocsvm,mirsky2018kitsune,ruff2018deep}. Our protocol follows this operational perspective through independent benign calibration and explicit reporting of observed test benign FPR.

From an applied-security perspective, this distinction is crucial: \emph{target FPR} is a calibration rule, while \emph{observed test benign FPR} is an operational outcome. Several anomaly-detection lines report ranking quality without tightly coupling threshold policy to independent benign calibration, which can overstate deployment readiness when traffic conditions shift \cite{mirsky2018kitsune,ruff2018deep,manocchio2024flowtransformer}. Our evaluation design therefore treats calibration and realized benign alert rate as first-class reporting items, alongside AUC and AP.

For explainability, attribution methods such as Integrated Gradients and SHAP are common in IDS and model auditing workflows \cite{sundararajan2017axiomatic,lundberg2017shap,rahman2025explaining}. However, attention mechanisms should not be treated as faithful explanations by default \cite{jain2019attention}. Our evaluation therefore combines attribution-guided masking with conservative interpretation, positioning auditability as empirical decision-trace support rather than proof of causal explanation.
